\chapter{Linear Regression and Least Squares}

\section{Linear Models and Projections}

Linear regression is one of the simplest and most fundamental models in machine learning. Despite its simplicity, it illustrates many key ideas that extend to more complex models.

\medskip

\noindent
\textbf{Model.}  
We consider functions of the form
\[
f(x) = \langle w, x \rangle + b,
\]
where $w \in \mathbb{R}^d$ and $b \in \mathbb{R}$.

\medskip

Given a dataset $\{(x_i, y_i)\}_{i=1}^n$, we aim to find parameters $(w,b)$ that best fit the data.

\medskip

\noindent
\textbf{Matrix formulation.}  
Let
\[
X =
\begin{bmatrix}
- & x_1^\top & - \\
- & x_2^\top & - \\
& \vdots & \\
- & x_n^\top & -
\end{bmatrix}
\in \mathbb{R}^{n \times d},
\quad
y =
\begin{bmatrix}
y_1 \\ y_2 \\ \vdots \\ y_n
\end{bmatrix}.
\]

For simplicity, we incorporate the bias into the parameter vector by augmenting $x$ with a constant component. Then the model becomes
\[
f(x) = \langle w, x \rangle,
\]
and the prediction vector is
\[
\hat{y} = X w.
\]

\medskip

\noindent
\textbf{Geometric interpretation.}  
The columns of $X$ span a subspace of $\mathbb{R}^n$. Linear regression seeks a vector $\hat{y}$ in this subspace that is closest to $y$.

\medskip

\noindent
Thus, linear regression can be viewed as an \emph{orthogonal projection} problem.

\section{Least Squares as Optimization}

A natural choice of loss is the squared error:
\[
\ell(f(x_i), y_i) = (f(x_i) - y_i)^2.
\]

\medskip

\noindent
\textbf{Objective.}
\[
\min_{w} \; \frac{1}{n} \|Xw - y\|_2^2.
\]

\medskip

\noindent
This is a convex quadratic optimization problem.

\medskip

\noindent
\textbf{Optimality condition.}  
Setting the gradient to zero yields the normal equations:
\[
X^\top X w = X^\top y.
\]

\medskip

\noindent
\textbf{Solution (when $X^\top X$ is invertible).}
\[
w^\star = (X^\top X)^{-1} X^\top y.
\]

\medskip

\noindent
\textbf{Projection interpretation.}  
The fitted values $\hat{y} = X w^\star$ satisfy
\[
y - \hat{y} \perp \mathrm{span}(X),
\]
meaning the residual is orthogonal to the column space of $X$.

\medskip

\noindent
\textbf{Insight.}  
The least squares solution is the best approximation of $y$ within the model class in terms of Euclidean distance.

\medskip

\noindent
\textbf{Overparameterized case.}  
If $X^\top X$ is not invertible (e.g., $d > n$), there are infinitely many solutions. In this case, one often selects the minimum-norm solution:
\[
w^\star = X^\top (X X^\top)^{-1} y.
\]

\medskip

\noindent
This solution is closely related to implicit regularization, as discussed in previous chapters.

\section{Conditioning and Numerical Stability}

While the closed-form solution provides insight, practical computation depends on numerical properties of the matrix $X^\top X$.

\medskip

\noindent
\textbf{Conditioning.}  
The condition number of $X^\top X$ measures sensitivity to perturbations:
\[
\kappa(X^\top X) = \frac{\lambda_{\max}}{\lambda_{\min}},
\]
where $\lambda_{\max}$ and $\lambda_{\min}$ are the largest and smallest eigenvalues.

\medskip

\noindent
\textbf{Ill-conditioning.}
\begin{itemize}
    \item Occurs when $\lambda_{\min}$ is very small
    \item Leads to unstable solutions
    \item Small changes in data can produce large changes in $w$
\end{itemize}

\medskip

\noindent
\textbf{Causes.}
\begin{itemize}
    \item Highly correlated features
    \item Redundant or nearly redundant inputs
\end{itemize}

\medskip

\noindent
\textbf{Practical consequences.}
\begin{itemize}
    \item Numerical instability in inversion
    \item Poor generalization
\end{itemize}

\medskip

\noindent
\textbf{Remedies.}
\begin{itemize}
    \item Feature scaling and normalization
    \item Removing redundant features
    \item Regularization (discussed next)
\end{itemize}

\medskip

\noindent
\textbf{Algorithmic note.}  
In practice, direct inversion is often avoided in favor of more stable methods such as QR decomposition or singular value decomposition (SVD).

\section{Ridge Regression}

Ridge regression introduces $\ell_2$ regularization into the least squares problem.

\medskip

\noindent
\textbf{Objective.}
\[
\min_{w} \; \|Xw - y\|_2^2 + \lambda \|w\|_2^2,
\quad \lambda > 0.
\]

\medskip

\noindent
\textbf{Solution.}
\[
w^\star = (X^\top X + \lambda I)^{-1} X^\top y.
\]

\medskip

\noindent
\textbf{Interpretation.}
\begin{itemize}
    \item Adds stability by shifting eigenvalues of $X^\top X$
    \item Reduces sensitivity to noise
    \item Controls the magnitude of parameters
\end{itemize}

\medskip

\noindent
\textbf{Geometric view.}  
Ridge regression shrinks the solution toward the origin, balancing data fit and parameter norm.

\medskip

\noindent
\textbf{Bias--variance tradeoff.}
\begin{itemize}
    \item Increases bias slightly
    \item Reduces variance significantly
\end{itemize}

\medskip

\noindent
\textbf{Connection to conditioning.}  
The matrix $X^\top X + \lambda I$ is always invertible for $\lambda > 0$, improving numerical stability.

\medskip

\noindent
\textbf{Bayesian interpretation.}  
Ridge regression corresponds to assuming a Gaussian prior on $w$.

\section{A Unifying Perspective}

Linear regression illustrates several central themes of machine learning:

\begin{itemize}
    \item Learning as projection onto a function space
    \item Optimization as solving structured problems
    \item Regularization as control of complexity and stability
\end{itemize}

\medskip

\noindent
\textbf{Three perspectives.}
\begin{itemize}
    \item \textbf{Geometric:} projection onto a subspace
    \item \textbf{Algebraic:} solving normal equations
    \item \textbf{Statistical:} estimating conditional expectation
\end{itemize}

\medskip

\noindent
These perspectives extend naturally to more complex models, including neural networks.

\section{Outlook}

This chapter introduced linear regression as a foundational model:
\begin{itemize}
    \item least squares formulation,
    \item geometric interpretation,
    \item numerical considerations,
    \item regularization via ridge regression.
\end{itemize}

\medskip

In the next chapter, we extend these ideas to classification through logistic regression, where probabilistic modeling plays a central role.

\medskip

\noindent
\textbf{Guiding principle.}  
Simple models often reveal deep structure. Understanding them provides insight into more complex learning systems.